\documentclass[11pt,a4paper]{article}
\usepackage[T1]{fontenc}
\usepackage[margin=26mm,headheight=14pt]{geometry}
\usepackage{amsmath,amssymb,amsthm,mathtools,bm}
\usepackage{newtxtext,newtxmath}
\usepackage{booktabs}
\usepackage{enumitem}
\usepackage{algorithm,algpseudocode}
\usepackage{float,placeins}
\usepackage{microtype}
\usepackage{hyperref}
\usepackage{fancyhdr}
\hypersetup{
  colorlinks=false,pdfborder={0 0 0},
  pdftitle={Higher-Order Combinatorial Structures under Low-Order Statistical Constraints:
            Transformations, Directed Constructions, and Realizability},
  pdfsubject={Combinatorics, discrete optimization, combinatorial design},
  pdfkeywords={higher-order structure, 2-switch, zero-sum trade, integer feasibility,
               residual compensation, reachability}
}
\setlength{\parindent}{1.2em}
\setlength{\parskip}{0pt}
\linespread{1.06}
\setlength{\emergencystretch}{2em}
\allowdisplaybreaks[1]
\numberwithin{equation}{section}
\newtheorem{proposition}{Proposition}[section]
\newtheorem{lemma}[proposition]{Lemma}
\newtheorem{corollary}[proposition]{Corollary}
\theoremstyle{definition}
\newtheorem{definition}[proposition]{Definition}
\floatname{algorithm}{Algorithm}
\numberwithin{algorithm}{section}
\pagestyle{fancy}
\fancyhf{}
\fancyhead[L]{\small Higher-Order Combinatorial Structures under Low-Order Constraints}
\fancyhead[R]{\small Sections 1--6}
\fancyfoot[C]{\thepage}
\renewcommand{\headrulewidth}{0.35pt}
\title{\bfseries Higher-Order Combinatorial Structures under\\[0.15em]
Low-Order Statistical Constraints\\[0.7em]
\large Transformations, Directed Constructions, and Realizability}
\author{}
\date{}
\begin{document}
\maketitle
\vspace{-3.4em}

\section{Introduction}\label{sec:intro}

Consider two combinatorial systems on six elements:
\begin{equation}
\begin{aligned}
X_A&=\bigl\{\{1,2,3\},\{1,4,5\},\{2,4,6\},\{3,5,6\}\bigr\},\\
X_B&=\bigl\{\{1,2,4\},\{1,3,5\},\{2,3,6\},\{4,5,6\}\bigr\}.
\end{aligned}
\label{eq:intro-example}
\end{equation}
Each system consists of four $3$-element blocks, and each element occurs exactly twice;
in the two systems every element pair also co-occurs the same number of times.
The triples themselves, however, are different.
For instance, $\{1,2,3\}$ occurs once in $X_A$ and does not occur in $X_B$,
whereas $\{1,2,4\}$ exhibits the opposite behavior.

\begin{table}[htbp]
\centering
\caption{Motivating example: identical low-order statistics, different third-order structure}
\label{tab:intro-example}
\begin{tabular}{cccc}
\toprule
Record & $X_A$ & $X_B$ & Block size\\
\midrule
1 & $\{1,2,3\}$ & $\{1,2,4\}$ & 3\\
2 & $\{1,4,5\}$ & $\{1,3,5\}$ & 3\\
3 & $\{2,4,6\}$ & $\{2,3,6\}$ & 3\\
4 & $\{3,5,6\}$ & $\{4,5,6\}$ & 3\\
\bottomrule
\end{tabular}
\end{table}

This simple construction reveals a fact worth investigating:
\emph{once the element frequencies and the pairwise co-occurrence frequencies are fixed,
the higher-order combinatorial structure may still retain freedom of choice.}
Studying this freedom requires answering three progressive questions at once:
how to change the higher-order structure while the low-order statistics remain unchanged;
how to realize a prescribed higher-order structure deliberately;
and which targets are realizable under finite resource constraints and reachable
through a given set of local operations.

The question arises wherever finite sets are used to express multivariate joint events.
An incidence matrix describes the membership relation between elements and blocks,
while first-order marginals and second-order co-occurrences provide a compact,
computable statistical summary.
If a combinatorial system is reconstructed from pairwise relations alone,
the structural information contributed by three or more elements acting jointly
remains undetermined.
Low-order statistical fitting therefore addresses \emph{how to match known statistics},
whereas higher-order structure construction further addresses
\emph{how to arrange the combinations that those statistics allow}.

We take a multiset of blocks as the basic object and use quota-preserving local switches
to connect statistical fitting with structural control.
A feasible initialization followed by residual-guided iteration first produces a reference
system with a definite fitting quality;
composite switches that preserve the low-order counts then expose higher-order variability.
The switches are subsequently used for the directed induction of prescribed structures,
and a layer-by-layer construction is completed by freezing genuine blocks and subtracting
their contribution from the remaining statistical targets.
Finally, association graphs, integer feasibility models, and state transition relations
determine the boundary of what can be constructed.

The main contributions of this paper are grouped into three interconnected aspects:
\begin{enumerate}[leftmargin=2.2em,itemsep=0.25em]
\item \textbf{Constructive characterization of higher-order variability.}
We give a zero-sum compensating operation composed of two switches, prove that it preserves
element frequencies, block sizes, and all second-order counts, and establish a quantitative
measure of higher-order frequency differences.
\item \textbf{Constrained construction of prescribed higher-order structures.}
We introduce directed induction, complete-record freezing, and residual compensation,
prove a strict equivalence between the residual count loss and the loss of the complete
system, and thereby keep layer-by-layer structural control aligned with one single global
statistical objective.
\item \textbf{Layered decision procedure for realizability and reachability.}
Association support graphs generate structural candidates, an integer model tests
realizability under finite resources, and a state graph of operations describes which
structures can be reached from a given initial system.
\end{enumerate}

The remainder of the paper follows the same logic.
Section~\ref{sec:model} gives the unified mathematical model and the three tasks;
Section~\ref{sec:baseline} gives feasible initialization and residual-guided baseline generation;
Section~\ref{sec:variability} studies higher-order structural transformations within
low-order statistical equivalence classes;
Section~\ref{sec:directed} develops directed construction and residual compensation;
Section~\ref{sec:boundary} studies finite realizability and operational reachability.

\section{Mathematical Model and Problem Definition}\label{sec:model}

This section fixes the combinatorial system, the ideal and empirical low-order statistics,
the hard constraints, and the three research tasks.
All subsequent methods are developed within this model.

\subsection{Combinatorial systems and their incidence representation}\label{subsec:model-incidence}

Let $V=\{1,\ldots,n\}$ be the ground set of elements, and let
\begin{equation}
\mathcal B=(B_1,\ldots,B_M),\qquad \varnothing\neq B_o\subseteq V
\label{eq:blocks}
\end{equation}
be a system comprising $M$ blocks.
Every block consists of distinct elements; the same element set may occur as several
distinct records, so $\mathcal B$ is understood as a multiset.
Let $k_o=|B_o|$ be the size of block $o$ and let
$M_k=\#\{o:k_o=k\}$ be the number of blocks of size $k$, with
\begin{equation}
N=\sum_{o=1}^M k_o=\sum_{k=1}^n kM_k
\label{eq:total-incidences}
\end{equation}
denoting the total number of element occurrences.

Represent the system by a binary incidence matrix $X\in\{0,1\}^{M\times n}$ with
$X_{oi}=\mathbf 1\{i\in B_o\}$.
The number of occurrences of element $i$ is
\begin{equation}
c_i(X)=\sum_{o=1}^M X_{oi},\qquad
\sum_i X_{oi}=k_o,\qquad \sum_i c_i(X)=N.
\label{eq:incidence-counts}
\end{equation}
Throughout the paper we fix the record size sequence $\boldsymbol k=(k_1,\ldots,k_M)$
and the element quotas $\boldsymbol c=(c_1,\ldots,c_n)$;
when no confusion can arise we write $c_i(X)=c_i$.

For $i\neq j$, define the pairwise joint containment count
$t_{ij}(X)=\sum_oX_{oi}X_{oj}$, and record the number of singleton blocks by
$s_i(X)=\#\{o:B_o=\{i\}\}$.
For any nonempty subset $S\subseteq V$, define
\begin{equation}
t_S(X)=\sum_{o=1}^M\mathbf 1\{S\subseteq B_o\},\qquad
\mathsf H_h(X)=\bigl(t_S(X)\bigr)_{S\in\binom{V}{h}},\quad h\geq3.
\label{eq:high-order}
\end{equation}
Accordingly, $\mathcal S_h(X)=\{S\in\binom Vh:t_S(X)>0\}$ is the set of realized
$h$-th order structures.
The quantity $t_S$ counts the complete records that \emph{contain} $S$ and therefore
applies naturally to blocks with more than $h$ elements.
We write $m_X(B)=\#\{o:B_o=B\}$ for the multiplicity of the complete type $B$.

Counting block by block yields the uniform higher-order total identity
\begin{equation}
T_h:=\sum_{S\in\binom Vh}t_S(X)=\sum_kM_k\binom kh,
\qquad \binom kh=0\ (k<h).
\label{eq:high-order-total}
\end{equation}
With the size distribution fixed, $T_h$ is constant.
When $T_h>0$, we use
\begin{equation}
d_h(X,Y)=\frac{\|\mathsf H_h(X)-\mathsf H_h(Y)\|_1}{2T_h}
\label{eq:high-order-distance}
\end{equation}
to measure the difference in $h$-th order frequencies between two systems;
its values lie in $[0,1]$.

\subsection{Ideal statistics, empirical statistics, and count targets}\label{subsec:model-statistics}

Let $\pi_B$ be an ideal probability distribution over nonempty blocks $B$, with
$\sum_B\pi_B=1$.
Write \mbox{$a_i=\Pr(i\in B)>0$},
$a_{ij}=\Pr(\{i,j\}\subseteq B)$, and
$\mu=\mathbb E_\pi |B|=\sum_i a_i$.
The ideal first-order marginals and second-order conditional association parameters
are defined as
\begin{equation}
p_i=\frac{a_i}{\mu},\qquad
q_{ij}=
\begin{cases}
a_{ij}/a_i,&i\ne j,\\
\pi_{\{i\}}/a_i,&i=j.
\end{cases}
\label{eq:ideal-statistics}
\end{equation}
Thus $P=(p_i)$ gives the ideal proportion of all element occurrences,
the off-diagonal entries of $Q=(q_{ij})$ give conditional co-occurrence probabilities,
and the diagonal entries give the conditional probability of \emph{forming a block on its own}.
This diagonal convention incorporates singleton records into the same counting model.
The population statistics satisfy
\begin{equation}
\sum_i p_i=1,\qquad p_iq_{ij}=p_jq_{ji}\quad(i\ne j).
\label{eq:ideal-balance}
\end{equation}

In a finite system the corresponding empirical quantities are
\begin{equation}
\widehat p_i(X)=\frac{c_i(X)}{N},\qquad
\widehat q_{ij}(X)=
\begin{cases}
t_{ij}(X)/c_i(X),&i\ne j,\\
s_i(X)/c_i(X),&i=j,
\end{cases}\quad(c_i(X)>0).
\label{eq:empirical-statistics}
\end{equation}
Hence $P,Q$ describe a given ideal law, whereas
$\widehat P(X),\widehat Q(X)$ are produced by a concrete finite system.
In particular,
$c_i(X)\widehat q_{ij}(X)=c_j(X)\widehat q_{ji}(X)=t_{ij}(X)$.

After fixing the total number of occurrences $N$, we transcribe the ideal statistics
into symmetric target counts
\begin{equation}
\tau_{ij}=Np_iq_{ij}=Np_jq_{ji}\quad(i<j),\qquad
\sigma_i=Np_iq_{ii}.
\label{eq:count-targets}
\end{equation}
Given nonnegative weights $w_{ij},v_i,\lambda$, define the count-based fitting loss
\begin{equation}
E(X)=\sum_{i<j}w_{ij}\bigl(t_{ij}(X)-\tau_{ij}\bigr)^2+
\lambda\sum_i v_i\bigl(s_i(X)-\sigma_i\bigr)^2.
\label{eq:global-loss}
\end{equation}
This objective always acts on the raw counts of the complete system and serves as the
uniform evaluation criterion for the subsequent search and for validation.

\subsection{Hard constraints, reference systems, and equivalence classes}\label{subsec:model-feasible}

Let $\mathcal Z\subseteq\binom V2$ be a set of structurally forbidden element pairs.
The feasible region satisfying the quotas, the per-record sizes, and the forbidden-pair
relations is defined as
\begin{equation}
\begin{aligned}
\Omega:=\Omega(\boldsymbol c,\boldsymbol k;\mathcal Z)
=\bigl\{X\in\{0,1\}^{M\times n}:&
\ \sum_iX_{oi}=k_o\ \ \forall o,\\
&\ \sum_oX_{oi}=c_i\ \ \forall i,\\
&\ X_{oi}+X_{oj}\le 1\
\ \ \forall o,\{i,j\}\in\mathcal Z\bigr\}.
\end{aligned}
\label{eq:omega}
\end{equation}
For a given tolerance $\varepsilon\ge0$, we write
$\Omega_\varepsilon=\{X\in\Omega:E(X)\le\varepsilon\}$.

Let $X_0\in\Omega$ be a reference system.
Taking the second-order counts \emph{actually realized} by $X_0$ as the standard,
we define the statistical equivalence class
\begin{equation}
\mathcal E(X_0)=\bigl\{X\in\Omega:
t_{ij}(X)=t_{ij}(X_0)\ \forall i<j,\
s_i(X)=s_i(X_0)\ \forall i\bigr\}.
\label{eq:equivalence-class}
\end{equation}
Inside this set, $\widehat P$, $\widehat Q$, and $E$ are all fixed,
whereas $\mathsf H_h$ is the object we wish to change and control.

For directed construction, let
\begin{equation}
\mathcal G=\{(U_\nu,m_\nu,\bowtie_\nu)\}_{\nu=1}^{J},
\qquad |U_\nu|\ge3,\quad m_\nu\in\mathbb Z_{\ge0},\quad
\bowtie_\nu\in\{\ge,=\},
\label{eq:directed-goals}
\end{equation}
denote a target set requiring $t_{U_\nu}(X)\bowtie_\nu m_\nu$.
Let $F$ denote the complete block records already frozen and $R$ the records that remain
mutable, and write $X=F\sqcup R$.
The multiset union $\sqcup$ is taken over concrete record indices.
Set $\Omega(F)=\{X\in\Omega:X\text{ preserves all records of }F\}$.

For any group of records $Y$, let $c_i(Y)$, $M_k(Y)$, $t_{ij}(Y)$, and $s_i(Y)$ be the
corresponding counts, and define the additive statistic vector
\begin{equation}
\Phi(Y)=\bigl((c_i(Y))_i,(M_k(Y))_k,
(t_{ij}(Y))_{i<j},(s_i(Y))_i\bigr).
\label{eq:phi-definition}
\end{equation}
This notation uniformly represents the element quotas, the record capacities, and the
low-order counts consumed by the frozen records.

To study association support, define
\begin{equation}
L_Q=\bigl\{\{i,j\}:i<j,\ q_{ij}>0\bigr\},\qquad
G_Q=(V,L_Q),\qquad
G_{\rm allow}=\left(V,\binom V2\setminus\mathcal Z\right).
\label{eq:support-graphs}
\end{equation}
The positive-association coverage feasible region is
\begin{equation}
\Omega^+(L_Q)=\{X\in\Omega:t_{ij}(X)\ge1,\ \forall\{i,j\}\in L_Q\}.
\label{eq:positive-support-domain}
\end{equation}
Positive-association coverage requires every specified positive association to be realized
at least once; when the task requires the support graph to be preserved exactly, one
additionally imposes $t_{ij}(X)=0$ for $\{i,j\}\notin L_Q$.
For any graph $G$, let
$\mathcal K_h(G)=\{S\in\binom Vh:\binom S2\subseteq E(G)\}$ denote all $h$-th order
clique candidates.

\subsection{Local switches}\label{subsec:model-switch}

The basic operation exchanges $a$ and $b$ between two distinct records
$A=A_0\cup\{a\}$ and $B=B_0\cup\{b\}$, where $a\in A\setminus B$ and $b\in B\setminus A$:
\begin{equation}
(A,B)\ \xrightarrow{\;2\text{-switch}\;}\
(A_0\cup\{b\},\,B_0\cup\{a\}).
\label{eq:2switch}
\end{equation}
This operation preserves the size of every record and all element quotas.
Checking the forbidden-pair condition on the two modified records yields a legal $2$-switch.
Records from different size strata are allowed to participate in the same switch;
their local count changes, residual-guided search, and feasible initialization are all
given uniformly in Section~\ref{sec:baseline}.

\subsection{Three progressive research questions}\label{subsec:model-questions}

\noindent\textbf{Problem 1: structural variability.}
Given $X_0\in\Omega$, construct $Y\in\mathcal E(X_0)$ such that
$\mathsf H_h(Y)\neq\mathsf H_h(X_0)$ for some $h\ge3$, and quantify the structural
difference by $d_h$.

\smallskip
\noindent\textbf{Problem 2: construction of a prescribed structure.}
Given $\mathcal G$ and a tolerance $\varepsilon$, find a system in $\Omega_\varepsilon$
that satisfies all targets.
In exact-preservation mode, the system is further required to belong to $\mathcal E(X_0)$.
The algorithm records realized targets and remaining resources in the form $F\sqcup R$.

\smallskip
\noindent\textbf{Problem 3: boundary decision.}
Given association support requirements $L_Q$, finite resources, and an allowed operation
set $\mathcal A$, determine which $h$-th order structures are realizable in
$\Omega^+(L_Q)$, and compare these structures with those reachable from $X_0$ through
$\mathcal A$.

\section{Feasible Generation under Low-Order Statistical Constraints}\label{sec:baseline}

Before studying higher-order structures we first need a reference system $X_0$ with a
definite low-order statistical quality.
This section constructs it: integer quotas are set from the ideal marginals, an initial
configuration satisfying the block capacities and the forbidden-pair conditions is found,
and legal $2$-switches are then iterated under the guidance of second-order count residuals.
The search preserves the hard constraints as invariants throughout and measures fitting
quality by the global count loss of \eqref{eq:global-loss}.

\subsection{Integer resource allocation}\label{subsec:baseline-budget}

Given the number of blocks $M$ and their size distribution $(M_k)$, expand the strata into
a sequence $\boldsymbol k=(k_1,\ldots,k_M)$; the total number of element occurrences is
$N=\sum_k kM_k$.
For ideal marginals $P=(p_i)$, the corresponding real-valued occurrence counts are $Np_i$.
In choosing integer quotas $\boldsymbol c$, one may solve
\begin{equation}
\begin{aligned}
\min_{\boldsymbol c\in\mathbb Z_{\ge0}^{n}}\quad&
\sum_{i=1}^{n}|c_i-Np_i|,\\
\mathrm{s.t.}\quad&
\sum_{i=1}^n c_i=N,\qquad 0\le c_i\le M\quad(i\in V).
\end{aligned}
\label{eq:baseline-margins}
\end{equation}
When the capacities allow, one may also round by the largest remainder method.
For elements whose empirical conditional association rows are to be fully estimated,
choose $c_i>0$.
Once the quotas are fixed, the first-order statistics of every legal system are fixed at
\begin{equation}
\widehat p_i(X)=\frac{c_i}{N},\qquad
\left|\widehat p_i(X)-p_i\right|
=\frac{|c_i-Np_i|}{N}.
\label{eq:baseline-margin-error}
\end{equation}
Integer rounding of the quotas therefore controls the first-order error, and the subsequent
iterations concentrate on adjusting second-order associations.

Integer quotas and record capacities must first satisfy
$\sum_i c_i=\sum_o k_o$, $c_i\le M$, and $k_o\le n$.
More precise feasibility is determined together with the incidence-matrix initialization
of the next step.

\subsection{Feasible initialization: maximum flow and integer programming}\label{subsec:baseline-init}

The initialization solves the following incidence-matrix problem:
\begin{equation}
\begin{aligned}
X^{(0)}\in\{0,1\}^{M\times n},\qquad&
\sum_iX^{(0)}_{oi}=k_o\quad(o=1,\ldots,M),\\
&\sum_oX^{(0)}_{oi}=c_i\quad(i\in V),\\
&X^{(0)}_{oi}+X^{(0)}_{oj}\le1
\quad(o=1,\ldots,M,\ \{i,j\}\in\mathcal Z).
\end{aligned}
\label{eq:baseline-initial-constraints}
\end{equation}
Initialization first establishes a strictly legal configuration;
second-order statistics are optimized afterwards by switch iteration.

When $\mathcal Z=\varnothing$, build a four-layer flow network:
source $s$, record nodes $o$, element nodes $i$, and sink $t$.
The capacities of the three classes of directed edges are
\begin{equation}
s\to o:\ k_o,\qquad o\to i:\ 1,\qquad i\to t:\ c_i.
\label{eq:baseline-network}
\end{equation}
An integral maximum flow of value $N$ saturates the capacity of every record node and
realizes the prescribed number of occurrences of each element.
Taking $X^{(0)}_{oi}$ as the flow on edge $o\to i$ yields the binary matrix of
\eqref{eq:baseline-initial-constraints}.

When $\mathcal Z\ne\varnothing$, add the per-record linear constraint associated with each
forbidden edge to the row-sum and column-sum equations, which directly yields a binary
integer program.
For instances with forbidden pairs one may also add predetermined record-type restrictions
to the feasibility solver and let the solver return a complete system satisfying the hard
constraints.

After $X^{(0)}$ is built, a single pass over all records and their internal element pairs
yields the initial $t_{ij}(X^{(0)})$ and $s_i(X^{(0)})$.
The cost of direct counting is
\begin{equation}
O\!\left(\sum_{o=1}^M k_o^2\right).
\label{eq:baseline-count-cost}
\end{equation}
The target counts $\tau_{ij},\sigma_i$ follow from \eqref{eq:count-targets}, and the initial
error $E(X^{(0)})$ follows from \eqref{eq:global-loss}.
The generation problem is thereby reduced to finding a system of lower error inside the
fixed feasible region $\Omega$.

\subsection{Cross-stratum switches and the search space}\label{subsec:baseline-switch}

Using the $2$-switch defined in Section~\ref{subsec:model-switch}, exchange $a$ and $b$
between two records $A=A_0\cup\{a\}$ and $B=B_0\cup\{b\}$.
The legality check involves only the newly inserted elements:
\begin{equation}
\{b,z\}\notin\mathcal Z\quad(\forall z\in A_0),\qquad
\{a,z\}\notin\mathcal Z\quad(\forall z\in B_0).
\label{eq:baseline-swap-legal}
\end{equation}
Every legal switch preserves the size of each record and all element quotas, so the search
state always lies in $\Omega$.

To enlarge the set of adjustable statistical dimensions, the two records are allowed to
come from different size strata.
For example,
\begin{equation}
\{\mathbf1\},\ \{\mathbf2,3\}
\ \longrightarrow\
\{\mathbf3\},\ \{\mathbf2,1\}.
\label{eq:baseline-cross-stratum}
\end{equation}
Both record sizes and the quotas of all three elements remain unchanged, while the element
pair $\{2,3\}$ is replaced by $\{1,2\}$, and the element occurring in the singleton record
changes from $1$ to $3$.

Cross-stratum switches have a definite statistical effect.
Define the number of occurrences of element $i$ in blocks of size $k$ as
\begin{equation}
d_{ik}(X)=\#\{o:|B_o|=k,\ i\in B_o\}.
\label{eq:baseline-layer-count}
\end{equation}
If switches are always confined to records of equal size, all $d_{ik}$ are invariants,
which simultaneously fixes
\begin{equation}
\widehat q_{ii}(X)=\frac{d_{i1}(X)}{c_i},\qquad
\sum_{j\ne i}\widehat q_{ij}(X)
=\frac{\sum_k(k-1)d_{ik}(X)}{c_i}.
\label{eq:baseline-layer-invariant}
\end{equation}
Cross-stratum switches, by contrast, allow the $d_{ik}$ to be redistributed among strata and
thus take part in the joint fitting of diagonal and off-diagonal terms.

\subsection{Second-order residuals and the exact gain of a single switch}\label{subsec:baseline-local-delta}

Let the current second-order count residuals be
\begin{equation}
r_{ij}(X)=t_{ij}(X)-\tau_{ij}\quad(i<j),\qquad
u_i(X)=s_i(X)-\sigma_i.
\label{eq:baseline-residuals}
\end{equation}
For a switch $X\to X'$, write
$\Delta t_{ij}=t_{ij}(X')-t_{ij}(X)$ and
$\Delta s_i=s_i(X')-s_i(X)$.
Counting record by record shows that, for $z\notin\{a,b\}$,
\begin{equation}
\begin{aligned}
\Delta t_{az}
&=\mathbf1\{z\in B_0\}-\mathbf1\{z\in A_0\},\\
\Delta t_{bz}
&=-\Delta t_{az},\qquad \Delta t_{ab}=0,\\
\Delta s_a
&=\mathbf1\{|B|=1\}-\mathbf1\{|A|=1\},
\qquad \Delta s_b=-\Delta s_a .
\end{aligned}
\label{eq:switch-local-change}
\end{equation}
All other element-pair and singleton counts remain unchanged.
Hence the second-order change of a single switch is confined to the count terms involving
$a,b$ and $A_0\triangle B_0$.

Using the difference identity for squared residuals, we obtain the \emph{exact} loss change
of a single switch
\begin{equation}
\boxed{
\begin{aligned}
\Delta E
={}&\sum_{z\in A_0\triangle B_0}
\sum_{x\in\{a,b\}}
w_{xz}\bigl(2r_{xz}\Delta t_{xz}+(\Delta t_{xz})^2\bigr)\\
&+\lambda\sum_{x\in\{a,b\}}
v_x\bigl(2u_x\Delta s_x+(\Delta s_x)^2\bigr).
\end{aligned}
}
\label{eq:switch-loss-delta}
\end{equation}
Here $w_{xz},r_{xz}$ are read symmetrically for unordered element pairs.
As long as the current residuals are directly accessible, the cost of evaluating a candidate
grows with $|A_0\triangle B_0|$ and does not require recomputing the full second-order count
matrix.

\paragraph{Computational example with two records.}
Continuing \eqref{eq:baseline-cross-stratum}, take targets $\tau_{12}=1$ and
$\tau_{13}=\tau_{23}=0$,
$\sigma_1=\sigma_2=0$ and $\sigma_3=1$,
set all weights to $1$, and let $\lambda=1$.
Before the switch the nonzero counts are $t_{23}=1$ and $s_1=1$;
after the switch they are $t_{12}=1$ and $s_3=1$.
Hence
\begin{equation}
E(X)=4,\qquad E(X')=0,\qquad\Delta E=-4.
\label{eq:baseline-swap-example}
\end{equation}
A single switch simultaneously corrects two element-pair counts and two singleton counts,
and \eqref{eq:switch-loss-delta} incorporates all these gains into the selection criterion.

\subsection{Residual-guided candidate generation and acceptance rules}\label{subsec:baseline-proposal}

Residuals not only evaluate the current system but also guide the next search step.
When $r_{ij}>0$ the pair $\{i,j\}$ co-occurs too often, and candidate generation
preferentially tries to separate the two elements;
when $r_{ij}<0$ it preferentially tries to make them co-occur.
A single switch usually affects several counts, so every candidate is re-evaluated by
\eqref{eq:switch-loss-delta} for its total gain.

Concretely, we use a mixed proposal.
With probability $\rho\in[0,1]$, select an element pair whose current error is prominent
according to $\lvert r_{ij}\rvert$ or a weighted version thereof, and look for a legal
switch among the records that contain this pair or are able to introduce it;
with probability $1-\rho$, sample two records and $a\in A\setminus B$,
$b\in B\setminus A$ at random.
Both proposal types share the same legality test and the same $\Delta E$ evaluation.

With pure discrete descent, only candidates with $\Delta E\le0$ are accepted.
To enlarge the search range, one may use a temperature $T_\ell>0$ that decreases with the
iteration and accept a legal switch with probability
\begin{equation}
\Pr(\mathrm{accept}\mid X,X')=
\begin{cases}
1,&\Delta E\le0,\\
\exp(-\Delta E/T_\ell),&\Delta E>0
\end{cases}
\label{eq:baseline-annealing}
\end{equation}
The algorithm also keeps the system $X_{\mathrm{best}}$ with the lowest error over the whole
search rather than directly returning the last search state.

\subsection{Error bounds and output validation}\label{subsec:baseline-properties}

Integer counting yields a direct second-order fitting baseline.
Let
\begin{equation}
I_{ij}=
\bigl[\max\{0,c_i+c_j-M\},\,\min\{c_i,c_j\}\bigr]\cap\mathbb Z.
\label{eq:baseline-pair-range}
\end{equation}

\begin{proposition}[Lower bound on the second-order loss under finite resources]
\label{prop:baseline-lower-bound}
For any $X\in\Omega$,
\begin{equation}
\begin{aligned}
E(X)\ge {}&
\sum_{i<j}w_{ij}\operatorname{dist}(\tau_{ij},I_{ij})^2\\
&+\lambda\sum_i v_i
\operatorname{dist}\bigl(\sigma_i,
\{0,\ldots,\min(c_i,M_1)\}\bigr)^2,
\end{aligned}
\label{eq:baseline-lower-bound}
\end{equation}
where $\operatorname{dist}(x,S)=\min_{y\in S}|x-y|$.
\end{proposition}
\begin{proof}
The number of joint occurrences of any two elements satisfies
$\max\{0,c_i+c_j-M\}\le t_{ij}(X)\le\min\{c_i,c_j\}$;
the number of singleton records satisfies
$0\le s_i(X)\le\min(c_i,M_1)$.
For each nonnegative weighted square, take the minimum over the realizable integer interval
and sum the results.
\end{proof}

\subsection{The complete generation algorithm and its interface to later sections}\label{subsec:baseline-algorithm}

The steps above are combined in Algorithm~\ref{alg:baseline}.
Integer quotas and feasible initialization establish a legal starting point;
residual guidance and random proposals supply the exploration directions;
the local increment formula provides exact evaluation;
the annealing rule decides the move from the current state, and the best-state record
determines the final output.

\begin{algorithm}[H]
\small
\caption{Feasible initialization and residual-guided second-order statistical fitting}
\label{alg:baseline}
\begin{algorithmic}[1]
\Require Ideal statistics $P,Q$; size sequence $\boldsymbol k$;
forbidden set $\mathcal Z$; weights; proposal ratio $\rho$;
temperature sequence $(T_\ell)$; iteration budget $L$; target threshold $\varepsilon$
\Ensure Best-seen reference system $X_0$ and its statistical error
\State $N\gets\sum_o k_o$
\State $\boldsymbol c\gets\Call{RoundMargins}{NP,\boldsymbol k}$
\State $X\gets\Call{FeasibleInit}{\boldsymbol c,\boldsymbol k,\mathcal Z}$
\State Compute target counts $\tau,\sigma$ and current counts $t,s$
\State $E\gets E(X)$; $X_{\mathrm{best}}\gets X$;
       $E_{\mathrm{best}}\gets E$
\For{$\ell=1,\ldots,L$}
  \State With probability $\rho$ choose a residual-guided proposal; otherwise use a random proposal
  \State Propose $(A,B,a,b)\gets\Call{ProposeSwap}{X,r,\rho}$
  \If{no candidate exists or the forbidden-pair check fails}
    \State \textbf{continue}
  \EndIf
  \State Compute $\Delta t,\Delta s$ from \eqref{eq:switch-local-change}
  \State Compute $\Delta E$ from \eqref{eq:switch-loss-delta}
  \If{$\Delta E\le0$ \textbf{or} $U<\exp(-\Delta E/T_\ell)$,
       where $U\sim\mathrm{Uniform}(0,1)$}
    \State Perform the switch; update $t,s,r,u,E$ locally
    \If{$E<E_{\mathrm{best}}$}
      \State $X_{\mathrm{best}}\gets X$;
             $E_{\mathrm{best}}\gets E$
    \EndIf
  \EndIf
  \If{$E_{\mathrm{best}}\le\varepsilon$}
    \State \textbf{break}
  \EndIf
\EndFor
\State $X_0\gets X_{\mathrm{best}}$
\State \Return $X_0,E_{\mathrm{best}},\widehat P(X_0),\widehat Q(X_0)$
\end{algorithmic}
\end{algorithm}

For a pure descent algorithm, the acceptance condition reduces to $\Delta E\le0$.
When a temperature sequence is used, let each $T_\ell>0$ and decrease its value gradually
according to the annealing schedule in use.
In an implementation it suffices to keep updating the few counts affected by each switch and
to maintain an inverted index from elements to records, so that residual-guided candidate
generation remains efficient.

The output $X_0$ preserves the given quotas, per-record sizes, and forbidden-pair relations;
its second-order fitting result is characterized by $E(X_0)$, $\widehat Q(X_0)$, and the
associated count residuals.
This step provides a common starting point for the three subsequent problems:

Section~\ref{sec:variability} fixes the low-order counts actually realized by $X_0$ and
explores higher-order structural changes within the statistical equivalence class;
Section~\ref{sec:directed} induces and freezes prescribed structures under the statistical
targets of the complete system;
Section~\ref{sec:boundary} further studies the structural boundary under finite resources and
operation rules.

\section{Higher-Order Structural Transformations under Fixed Low-Order Statistics}\label{sec:variability}

This section answers the first question: how can the higher-order structure be changed while
the complete second-order statistics remain unchanged?
The core idea is to organize two local switches into a composite operation whose low-order
count increments cancel.

\subsection{Describing a switch by its count increment}\label{subsec:variation-signature}

Let $X'$ be obtained from $X$ by one legal $2$-switch.
Combine the pairwise co-occurrence changes and the singleton changes into the symmetric matrix
\begin{equation}
D_{ij}(X,X')=
\begin{cases}
t_{ij}(X')-t_{ij}(X),&i\neq j,\\
s_i(X')-s_i(X),&i=j.
\end{cases}
\label{eq:low-order-signature}
\end{equation}
Since a switch preserves element quotas and per-record sizes, $D$ completely describes the
change in the remaining low-order statistics.
For $A_0,B_0,a,b$ in \eqref{eq:2switch}, let $\boldsymbol e_i$ be the $i$-th unit vector and
$\boldsymbol 1_C$ the indicator vector of a set $C$, and write
\begin{equation}
\boldsymbol r=\boldsymbol1_{B_0}-\boldsymbol1_{A_0},\qquad
\delta=\mathbf1\{|B|=1\}-\mathbf1\{|A|=1\}.
\label{eq:trade-signature}
\end{equation}
Then \eqref{eq:switch-local-change} can be written as
\begin{equation}
D=(\boldsymbol e_a-\boldsymbol e_b)\boldsymbol r^\top+
\boldsymbol r(\boldsymbol e_a-\boldsymbol e_b)^\top+
\delta\operatorname{Diag}(\boldsymbol e_a-\boldsymbol e_b).
\label{eq:signature-matrix}
\end{equation}
This expression compresses the low-order effect of a switch into $(a,b,\boldsymbol r,\delta)$.
For the same element pair $a,b$, two switches whose background difference vectors and
singleton changes are both opposite have low-order effects that cancel exactly.

\subsection{Zero-sum compensating trades on four blocks}\label{subsec:variation-trade}

On four records with pairwise distinct indices, take
\begin{equation}
\begin{aligned}
A&=A_0\cup\{a\},&B&=B_0\cup\{b\},\\
C&=C_0\cup\{a\},&D&=D_0\cup\{b\},
\end{aligned}
\label{eq:four-block-trade}
\end{equation}
where the two switches take place on $(A,B)$ and $(C,D)$ respectively.
Let $(\boldsymbol r_1,\delta_1)$ and $(\boldsymbol r_2,\delta_2)$ denote the signatures of
the two switches.

\begin{proposition}[Low-order preservation of a zero-sum compensating trade]
\label{prop:zero-sum-trade}
If the four records have pairwise distinct indices, both switches are legal, and
\begin{equation}
\boldsymbol r_1+\boldsymbol r_2=\boldsymbol0,\qquad
\delta_1+\delta_2=0,
\label{eq:zero-sum-condition}
\end{equation}
then the system $Y$ obtained after the composite switch belongs to $\mathcal E(X)$.
\end{proposition}
\begin{proof}
Each switch preserves all row sums and column sums.
By \eqref{eq:signature-matrix} and \eqref{eq:zero-sum-condition}, the low-order change
matrices of the two switches satisfy $D_1+D_2=0$.
Since the four records are distinct, the two operations can be carried out independently on
the same original system, so the composite operation preserves all $t_{ij}$ and $s_i$.
Combining the row- and column-sum invariants with legality yields $Y\in\mathcal E(X)$.
\end{proof}

We return to the system $X_A$ of \eqref{eq:intro-example}.
The first step exchanges $3$ and $4$ between $\{1,2,\mathbf3\}$ and $\{1,\mathbf4,5\}$,
whose background difference vector is
$\boldsymbol r_1=\boldsymbol e_5-\boldsymbol e_2$.
The second step exchanges the same pair of elements between $\{\mathbf3,5,6\}$ and
$\{2,\mathbf4,6\}$, whose background difference vector is
$\boldsymbol r_2=\boldsymbol e_2-\boldsymbol e_5$.
Both steps take place on $3$-element records, so $\delta_1=\delta_2=0$.
The composite operation yields exactly $X_B$ and satisfies
\begin{equation}
X_B\in\mathcal E(X_A),\qquad
\mathcal S_3(X_A)\cap\mathcal S_3(X_B)=\varnothing.
\label{eq:example-invariant}
\end{equation}
Each system has four triples,
$\|\mathsf H_3(X_A)-\mathsf H_3(X_B)\|_1=8$ and $T_3=4$, hence
\begin{equation}
d_3(X_A,X_B)=\frac{8}{2\times4}=1.
\label{eq:example-distance}
\end{equation}
This exhibits a concrete transformation that attains the maximum normalized third-order
difference while the complete second-order counts are conserved.

\subsection{Searching for higher-order structures within an equivalence class}\label{subsec:variation-search}

A zero-sum compensating trade can serve directly as a search step for exploring the
equivalence class.
Starting from $X_0$, the algorithm first enumerates or samples a first legal switch and
computes its signature $(a,b,\boldsymbol r,\delta)$;
it then looks among the unused records for a legal switch with signature
$(a,b,-\boldsymbol r,-\delta)$.
The outcomes of the two switches form a candidate $Y$, which is accepted when its
higher-order frequency vector changes, and the local indices and higher-order counts are
updated.

\begin{algorithm}[H]
\small
\caption{Higher-order structure exploration within a second-order statistical equivalence class}
\label{alg:variation}
\begin{algorithmic}[1]
\Require Legal reference system $X_0$; forbidden set $\mathcal Z$; search budget $L$
\Ensure Construction $X$ preserving the low-order statistics of the reference system
\State $X\gets X_0$
\For{$\ell=1,\ldots,L$}
  \State Select a first legal switch in $X$ and compute its signature $(a,b,\boldsymbol r,\delta)$
  \State Search the remaining records for a legal switch with signature $(a,b,-\boldsymbol r,-\delta)$
  \If{a compensating switch is found}
    \State Compute the system $Y$ after the composite switch on the four records
    \If{$\mathsf H_h(Y)\ne\mathsf H_h(X)$ for some $h\ge3$ of interest}
      \State $X\gets Y$; update the affected records and higher-order counts
    \EndIf
  \EndIf
\EndFor
\State \Return $X$
\end{algorithmic}
\end{algorithm}

The search can be directed toward different structural objectives.
If higher-order diversity is of interest, one may record $\lvert\mathcal S_h(X)\rvert$ and
retain several distinct higher-order frequency vectors;
if one wishes to move away from the initial system, one may preferentially accept candidates
that increase $d_h(X,X_0)$.
In an implementation, switch signatures can be matched by hash tables or inverted indices,
so that the compensation search concentrates on a small number of relevant records.

The construction above reveals a definite structural operation: by designing the change in
the low-order counts as a zero-sum combination, one can move along the statistical
equivalence class to change the higher-order structure.
The next section advances this structural transformation from \emph{exploring different
outcomes} to \emph{realizing a prescribed outcome}.

\section{Directed Construction of Prescribed Higher-Order Structures}\label{sec:directed}

Section~\ref{sec:variability} exhibited higher-order variability within a low-order
statistical equivalence class by means of zero-sum compensating trades.
This section further considers the problem of realizing a prescribed target set $\mathcal G$.
The basic procedure is first to \emph{induce} the target, choosing the records that actually
carry it as a temporary anchor group, then to \emph{freeze} these records and perform
\emph{residual compensation} on the free part, and finally to validate the complete
candidate, \emph{committing} it officially only if validation succeeds.
Intermediate states in a candidate construction need not satisfy the final statistical
conditions; an already committed complete system, by contrast, always satisfies the
corresponding conditions.

\subsection{Directed induction and target anchoring}\label{subsec:directed-induction}

Consider a target $U\subseteq V$ with $|U|\ge3$.
For a switch as in \eqref{eq:2switch}, the increment in the number of occurrences of $U$ is
\begin{equation}
\Delta_U=
\begin{cases}
\mathbf1\{U\setminus\{a\}\subseteq B_0\}
-\mathbf1\{U\setminus\{a\}\subseteq A_0\},
&a\in U,\ b\notin U,\\
\mathbf1\{U\setminus\{b\}\subseteq A_0\}
-\mathbf1\{U\setminus\{b\}\subseteq B_0\},
&b\in U,\ a\notin U,\\
0,&\text{otherwise}.
\end{cases}
\label{eq:directed-gain}
\end{equation}
Hence, to increase the number of structures $U$ containing $a$, one may preferentially look
for records that already contain $U\setminus\{a\}$ and move $a$ into such a record by a legal
switch.
Equality-type targets may also require decreasing an existing $U$, and the same increment
formula can be used to evaluate candidates.

For at-least targets, let $\eta_\nu>0$ and define the coverage score
\begin{equation}
S_{\mathcal G}(X)=
\sum_{\nu:\bowtie_\nu=\ge}
\eta_\nu\min\{m_\nu,t_{U_\nu}(X)\}.
\label{eq:directed-score}
\end{equation}
In the candidate-generation phase, one may rank switches by
\begin{equation}
J(X,Y)=S_{\mathcal G}(Y)-S_{\mathcal G}(X)
-\beta\bigl(E(Y)-E(X)\bigr),\qquad \beta\ge0,
\label{eq:directed-priority}
\end{equation}
and for equality-type targets the search is additionally guided by the decrease of
$|t_{U_\nu}(X)-m_\nu|$.
The score is used only to select candidates; it does not replace the final feasibility and
target validation.

Keeping the current frozen group $F$ fixed, only the free records are modified, yielding a
provisional candidate.
From the complete records that actually contain the target $U$ in the candidate, choose an
\emph{anchor group} $G$; records are distinguished by index, so that even identical element
sets count as different records.
By additivity of the counts,
\begin{equation}
t_U(F\sqcup R)=t_U(F)+t_U(R)\ge t_U(F).
\label{eq:anchor-guarantee}
\end{equation}
Hence, for $t_U\ge m$, as long as $F\sqcup G$ already contains at least $m$ records
containing $U$, this \emph{lower bound} is guaranteed permanently by never modifying these
frozen records.
For $t_U=m$, even if $m$ records have been frozen, the free records may still produce $U$
again, so equality must be checked at this and at every later official commit.
In particular, when $m=0$ no carrying record needs to be frozen for this target, and one may
take $G=\varnothing$.

We distinguish two final validation modes:
\emph{exact-preservation mode} requires the committed complete system $X'$ to satisfy
$X'\in\mathcal E(X_0)$;
\emph{tolerance-controlled mode} requires $X'\in\Omega$ and $E(X')\le\varepsilon$.
Both modes allow general legal $2$-switches on candidates that have not yet been committed,
temporarily changing second-order counts and compensating afterwards on the free part.
The zero-sum compensating trades of Section~\ref{sec:variability} remain directly usable as a
candidate-generation method, but there is no requirement that every internal candidate switch
preserve the second-order counts individually.

\subsection{Frozen resources and residual compensation}\label{subsec:directed-residual}

After the $r$-th successful commit, the cumulative frozen records are $F^{(r)}$, the free
records are $R_r$, and the complete system is written
\begin{equation}
X_r=F^{(r)}\sqcup R_r,\qquad F^{(0)}=\varnothing.
\label{eq:recursive-split}
\end{equation}
For the statistic vector of Section~\ref{sec:model},
\begin{equation}
\Phi(F\sqcup R)=\Phi(F)+\Phi(R).
\label{eq:stat-add}
\end{equation}
We combine the raw integer resources with the ideal low-order targets into
\begin{equation}
b^\star=\bigl(\boldsymbol c,(M_k)_k,
(\tau_{ij})_{i<j},(\sigma_i)_i\bigr).
\label{eq:original-budget}
\end{equation}
After freezing $F^{(r)}$, the remaining resources and fitting targets of the free system are
\begin{equation}
\boxed{\quad b_r^\star=b^\star-\Phi(F^{(r)}).\quad}
\label{eq:budget-subtraction}
\end{equation}
that is,
\begin{equation}
\begin{aligned}
d_i^{(r)}&=c_i-c_i(F^{(r)}),&
M_k^{(r)}&=M_k-M_k(F^{(r)}),\\
\tau_{ij}^{(r)}&=\tau_{ij}-t_{ij}(F^{(r)}),&
\sigma_i^{(r)}&=\sigma_i-s_i(F^{(r)}).
\end{aligned}
\label{eq:residual-targets}
\end{equation}
Here $d_i^{(r)},M_k^{(r)}$ are the integer resource constraints that the free system must
satisfy, while $\tau_{ij}^{(r)},\sigma_i^{(r)}$ are the fitting references in the loss
function, which need not be nonnegative integers.
After adding an anchor group $G$, one only needs to update
$b_{r+1}^\star=b_r^\star-\Phi(G)$.

With the original weights, define the residual loss
\begin{equation}
E_r(R)=
\sum_{i<j}w_{ij}\bigl(t_{ij}(R)-\tau_{ij}^{(r)}\bigr)^2+
\lambda\sum_i v_i\bigl(s_i(R)-\sigma_i^{(r)}\bigr)^2.
\label{eq:residual-loss}
\end{equation}

\begin{proposition}[Equivalence of the residual loss and the loss of the complete system]
\label{prop:residual-equivalence}
Fix any complete group of records $F^{(r)}$.
For any free record system $R$,
\begin{equation}
\boxed{\ E(F^{(r)}\sqcup R)=E_r(R).\ }
\label{eq:loss-equivalence}
\end{equation}
\end{proposition}
\begin{proof}
Additivity of the counts gives
\[
t_{ij}(F^{(r)}\sqcup R)-\tau_{ij}
=t_{ij}(R)-\bigl(\tau_{ij}-t_{ij}(F^{(r)})\bigr)
=t_{ij}(R)-\tau_{ij}^{(r)}.
\]
The singleton counts likewise satisfy
$s_i(F^{(r)}\sqcup R)-\sigma_i=s_i(R)-\sigma_i^{(r)}$.
Substituting the two identities into \eqref{eq:global-loss} gives the claim.
\end{proof}

This proposition shows that, as long as the frozen group is unchanged, optimizing $E_r$ on
the free system is optimizing the loss of the same complete system.
One may therefore reuse the legal switches, the local residual updates, and the search
methods of Section~\ref{sec:baseline} without redefining the global evaluation criterion.

For \emph{exact-preservation mode}, one must additionally distinguish the ideal targets from
the counts actually realized by the reference system.
After fixing $F^{(r)}$, the complete system belongs to $\mathcal E(X_0)$ if and only if its
hard constraints hold and the free part satisfies
\begin{equation}
\begin{aligned}
t_{ij}(R)&=t_{ij}(X_0)-t_{ij}(F^{(r)})&&(\forall i<j),\\
s_i(R)&=s_i(X_0)-s_i(F^{(r)})&&(\forall i).
\end{aligned}
\label{eq:exact-residual-targets}
\end{equation}
Exact-preservation mode therefore uses \eqref{eq:exact-residual-targets} as the final count
condition; unless the ideal counts happen to coincide with the actual counts of $X_0$,
checking $E_r(R)$ alone is not sufficient to decide membership in $\mathcal E(X_0)$.

\subsection{Recursive directed construction and correctness}\label{subsec:directed-algorithm}

Let $X_0\in\Omega$ and, in tolerance-controlled mode, additionally $E(X_0)\le\varepsilon$.
Let $\mathcal C$ denote the set of \emph{accepted target commitments}:
every $g=(U,m,\bowtie)\in\mathcal C$ requires each subsequently committed complete system to
continue satisfying $t_U\bowtie m$.
Unlike the frozen records $F$, $\mathcal C$ records \emph{logical conditions} that must be
checked continuously; a single frozen record cannot by itself guarantee an equality-type
target.

In each round, choose a pending target from $\mathcal G\setminus\mathcal C$, perform directed
induction on the free system, obtain a candidate $\widetilde R$ and a temporary anchor group
$G$, and set
\begin{equation}
\widetilde F=F\sqcup G,\qquad
\widetilde R_0=\widetilde R\setminus G.
\label{eq:temporary-freeze}
\end{equation}
Perform residual compensation on $\widetilde R_0$ according to
$b^\star-\Phi(\widetilde F)$, obtaining $R'$ and the complete candidate
$X'=\widetilde F\sqcup R'$.
Candidate generation allows the case of no switch (the target already occurs) and the case
$G=\varnothing$ (for instance, eliminating all occurrences of an equality-type target);
however, every \emph{successful commit} must turn at least one not-yet-committed target into
a new commitment.

Formal validation simultaneously checks the hard constraints of the complete system, the
previously frozen records, all existing target commitments, the new target of the current
round, and the statistical condition of the chosen mode:
\begin{equation}
\operatorname{Valid}_{\kappa}(X')=
\begin{cases}
\bigl[X'\in\mathcal E(X_0)\bigr],&
\kappa=\mathrm{exact},\\
\bigl[X'\in\Omega,\ E(X')\le\varepsilon\bigr],&
\kappa=\mathrm{tol}.
\end{cases}
\label{eq:directed-validation}
\end{equation}
Here the square brackets denote that the condition holds.
The candidate is committed officially only after all these checks pass.


\begin{proposition}[Validity of the recursive directed construction]
\label{prop:directed-validity}
If the initial system satisfies the prescribed hard constraints (and, in tolerance mode, the
initial error threshold), and before every commit the algorithm verifies by
\eqref{eq:directed-validation} the statistical condition of the complete system, the frozen
records, and all existing and new target commitments, then every complete system after a
successful commit satisfies the constraints of the chosen mode, and the targets accepted
earlier remain valid.
If the algorithm returns $\mathrm{SUCCESS}$, the output system satisfies all of $\mathcal G$.
\end{proposition}
\begin{proof}
Induct on the number of successful commits.
The initial system satisfies the corresponding constraints, and the initial commitment set
consists of the targets that hold in $X_0$, so the claim holds.
Assume that the current system and its commitment set satisfy the conclusion.
The next commit occurs only after the complete candidate has passed validation of the hard
constraints, the statistical condition, the frozen records, and $\mathcal C\cup\{g\}$;
hence the existing commitments are preserved and the new commitment $g$ holds.
The updated commitment set contains only the targets actually satisfied by the new system,
so the induction is complete.
When $\mathcal C=\mathcal G$, all targets hold.
\end{proof}

\begin{algorithm}[H]
\small
\caption{Recursive directed construction of prescribed higher-order structures}
\label{alg:directed}
\begin{algorithmic}[1]
\Require $X_0\in\Omega$; target set $\mathcal G$; mode
$\kappa\in\{\mathrm{exact},\mathrm{tol}\}$; tolerance $\varepsilon$;
per-layer budget $L$; maximum number of layers $r_{\max}$
\Ensure Validated system $X$, frozen records $F$, commitment set $\mathcal C$, and status
\State $X\gets X_0$; $F\gets\varnothing$;
$\mathcal C\gets\{g\in\mathcal G:g\text{ holds in }X_0\}$
\For{$r=1,\ldots,r_{\max}$}
  \If{$\mathcal C=\mathcal G$}
    \State \textbf{break}
  \EndIf
  \State $R\gets X\setminus F$; $\mathrm{accepted}\gets\mathrm{false}$
  \For{$\ell=1,\ldots,L$}
    \State Choose $g\in\mathcal G\setminus\mathcal C$;
           search for $(\widetilde R,G)$ in $R$
    \If{no candidate satisfying the hard constraints is found}
      \State \textbf{continue}
    \EndIf
    \State $\widetilde F\gets F\sqcup G$;
           $\widetilde R_0\gets\widetilde R\setminus G$
    \State $b'\gets b^\star-\Phi(\widetilde F)$
    \State $R'\gets\Call{ResidualRefit}{\widetilde R_0,b',\kappa,X_0}$
    \If{$R'=\bot$}
      \State \textbf{continue}
    \EndIf
    \State $X'\gets\widetilde F\sqcup R'$
    \If{$\operatorname{Valid}_{\kappa}(X')$
       and $X'$ preserves $F$ and satisfies $\mathcal C\cup\{g\}$}
      \State $X\gets X'$; $F\gets\widetilde F$
      \State $\mathcal C\gets\{h\in\mathcal G:h\text{ holds in }X\}$
      \State $\mathrm{accepted}\gets\mathrm{true}$ \Comment{official commit}
      \State \textbf{break}
    \EndIf
  \EndFor
  \If{$\mathrm{accepted}=\mathrm{false}$}
    \State \textbf{break}
  \EndIf
\EndFor
\If{$\mathcal C=\mathcal G$}
  \State $\mathrm{status}\gets\mathrm{SUCCESS}$
\ElsIf{$\mathcal C\ne\varnothing$}
  \State $\mathrm{status}\gets\mathrm{PARTIAL}$
\Else
  \State $\mathrm{status}\gets\mathrm{FAILURE}$
\EndIf
\State \Return $(X,F,\mathcal C,\mathrm{status})$
\end{algorithmic}
\end{algorithm}

Every successful commit adds at least one new target commitment, so the number of successful
commits does not exceed $|\mathcal G|$; the set of frozen records only grows or stays the
same.
$\mathrm{PARTIAL}$ and $\mathrm{FAILURE}$ mean only that the given operations, budget, and
search strategy did not complete all targets; they do \emph{not} constitute a proof that the
targets are mathematically unrealizable.

\subsection{A complete construction example: induction, freezing, and residual compensation}\label{subsec:directed-example}

Continue with the six-element system of \eqref{eq:intro-example}.
Below we abbreviate $\{1,2,3\}$ by $123$, and take
\begin{equation}
X_0=\{123,145,246,356\},\qquad
t_{124}\ge1,\quad t_{456}\ge1.
\label{eq:directed-example-initial}
\end{equation}
Initially neither target occurs.
There are no forbidden pairs, all four records have size $3$, and every element has quota
$2$.
Take $\tau_{ij}=t_{ij}(X_0)$ and $\sigma_i=s_i(X_0)=0$, set all second-order weights to
$1$, and use exact-preservation mode.
Since all records have size $3$, the singleton terms vanish identically;
hence $E(X_0)=0$ in this example, and validation requires the final system to restore all
original second-order counts.

\paragraph{Step 1: directed induction.}
Exchange elements $3$ and $4$ in records $123$ and $145$:
\begin{equation}
123,\ 145\longrightarrow124,\ 135,\qquad
\widetilde X=\{124,135,246,356\}.
\label{eq:directed-example-induction}
\end{equation}
Now $t_{124}(\widetilde X)=1$, but the second-order counts temporarily deviate from the
reference system, so the candidate cannot yet be committed in exact-preservation mode.

\paragraph{Step 2: temporary freezing and resource subtraction.}
Choose the newly produced record $124$ as the anchor record and fix it only within this
candidate branch:
\begin{equation}
\begin{aligned}
\widetilde F&=\{124\},&
\widetilde R_0&=\{135,246,356\},\\
\boldsymbol d^{(1)}&=(1,1,2,1,2,2),&
M_3^{(1)}&=3.
\end{aligned}
\label{eq:level-one-example}
\end{equation}
By \eqref{eq:residual-targets}, the target counts of the free system on the four deviating
element pairs are
\[
(\tau_{23}^{(1)},\tau_{24}^{(1)},
  \tau_{35}^{(1)},\tau_{45}^{(1)})=(1,0,1,1),
\]
whereas the corresponding actual counts of $\widetilde R_0$ are $(0,1,2,0)$.
Hence $E_1(\widetilde R_0)=E(\widetilde X)=4$.
This is only a temporary frozen state that has \emph{not} been committed; it does not mean
that a system violating the statistical condition has been accepted.

\paragraph{Step 3: residual compensation of the free part.}
Without touching the anchor record $124$, exchange elements $4$ and $3$ between the two free
records $246$ and $356$:
\begin{equation}
246,\ 356\longrightarrow236,\ 456,\qquad
R'=\{135,236,456\}.
\label{eq:directed-example-compensation}
\end{equation}
The residual counts of the free part are now restored, $E_1(R')=0$, and the compensation has
additionally produced the second target $456$.

\begin{table}[htbp]
\centering
\caption{Stepwise change of the second-order counts in the construction example}
\label{tab:directed-example-counts}
\begin{tabular}{cccc}
\toprule
Element pair & $\Delta$ at induction & $\Delta$ at compensation & Total $\Delta$\\
\midrule
$\{2,3\}$ & $-1$ & $+1$ & $0$\\
$\{2,4\}$ & $+1$ & $-1$ & $0$\\
$\{3,5\}$ & $+1$ & $-1$ & $0$\\
$\{4,5\}$ & $-1$ & $+1$ & $0$\\
Other pairs & $0$ & $0$ & $0$\\
\bottomrule
\end{tabular}
\end{table}

\paragraph{Step 4: overall validation and official commit.}
The complete candidate is
\begin{equation}
X'=\{124,135,236,456\},\qquad
E(X_0)=0,\quad E(\widetilde X)=4,\quad E(X')=0.
\label{eq:directed-example-acceptance}
\end{equation}
Table~\ref{tab:directed-example-counts} shows that all second-order counts have been
restored; at the same time the size of every record and the quota of every element have
remained unchanged, and $t_{124}(X')=t_{456}(X')=1$.
Hence $X'\in\mathcal E(X_0)$, so the candidate can be committed officially and both targets
enter the commitment set.
If one wants to fix the actual carrier of the second target as well, one may add $456$ to the
frozen group after this validation, obtaining
\begin{equation}
F^{(2)}=\{124,456\},\qquad R_2=\{135,236\},\qquad
\boldsymbol d^{(2)}=(1,1,2,0,1,1),\quad T_3(R_2)=2.
\label{eq:level-two-example}
\end{equation}
This additional freezing does not change the validated complete system and requires no
further switch.

What this example demonstrates is not a new zero-sum trade theorem:
its two local switches together are still a staged implementation of the four-record zero-sum
compensating trade of Section~\ref{sec:variability}.
What needs to be shown here is that the procedure of \emph{first inducing, then temporarily
anchoring, then compensating on the free part, and finally validating the whole} does
preserve the prescribed targets while repairing a nonzero residual.

\subsection{The effect of freezing on the construction space}\label{subsec:directed-budget}

Freezing consumes records and element quotas, and it also fixes the higher-order statistical
contribution of these records.
At layer $r$, the total $h$-th order containment count that the free system can carry is
\begin{equation}
T_h(R_r)=\sum_kM_k^{(r)}\binom kh,\qquad
T_h=T_h(F^{(r)})+T_h(R_r).
\label{eq:residual-high-budget}
\end{equation}

\begin{proposition}[Bound on higher-order variation under a common frozen group]
\label{prop:frozen-distance}
Let $X=F^{(r)}\sqcup R$ and
$Y=F^{(r)}\sqcup R'$,
and assume that the two free systems have the same size distribution.
For $T_h>0$,
\begin{equation}
\boxed{\
d_h(X,Y)\le\frac{T_h(R_r)}{T_h}
=1-\frac{T_h(F^{(r)})}{T_h}.\ }
\label{eq:frozen-distance}
\end{equation}
\end{proposition}
\begin{proof}
The higher-order frequencies of the common frozen group cancel in the difference, so
$\mathsf H_h(X)-\mathsf H_h(Y)
=\mathsf H_h(R)-\mathsf H_h(R')$.
Both free systems have nonnegative higher-order frequencies, and their component sums are
both $T_h(R_r)$; hence the $\ell_1$ distance is at most $2T_h(R_r)$.
Substituting into \eqref{eq:high-order-distance} gives the claim.
\end{proof}

If $F^{(r)}\subseteq F^{(r+1)}$, then $\Omega(F^{(r+1)})\subseteq\Omega(F^{(r)})$;
when these feasible regions are nonempty,
$\min_{X\in\Omega(F^{(r)})}E(X)$ is nondecreasing as freezing increases.
The bound above shows that the more higher-order counts are frozen in common, the smaller
the upper bound on the higher-order difference that two subsequent candidates can exhibit in
the complete system.
In the example above, $T_3(R)=3$ when $124$ is frozen and drops to $2$ after $456$ is also
frozen, which precisely reflects this contraction of the construction space.

\section{Realizability and Reachability of Higher-Order Structures}\label{sec:boundary}

The two preceding sections established methods for transforming and deliberately
constructing higher-order structures.
This section studies the structural boundary that these methods face.
To this end we organize the filtering of higher-order candidates into three layers:
\emph{the association graph supplies candidates, finite resources determine realizability,
and operation paths determine reachability.}

\subsection{Higher-order candidates from the pairwise association graph}\label{subsec:boundary-graph}

For a graph $G=(V,E_G)$, any $h$-th order candidate $S\in\mathcal K_h(G)$ satisfies
$\binom S2\subseteq E_G$;
that is, every two elements of $S$ have an allowed pairwise association.
If $N_G(i)$ denotes the set of neighbors of element $i$, then, for a given clique $S$, the
extension candidates are
\begin{equation}
C_G(S)=\left(\bigcap_{i\in S}N_G(i)\right)\setminus S.
\label{eq:clique-extension}
\end{equation}
One enumerates the cliques of every order by starting from the edges, iteratively adding the
elements of $C_G(S)$, and deduplicating by a fixed element order.

The two graphs introduced in Section~\ref{subsec:model-feasible} serve different construction
requirements.
When exact preservation of the association support is prescribed, candidates are enumerated
from $G_Q$;
when positive-association coverage is prescribed with additional co-occurrences allowed,
$G_{\rm allow}$ gives the complete range of allowed structures, and the cliques of $G_Q$ are
treated as reference candidates to be examined first.

For the motivating example $X_A$, the actual support graph contains the edges $\{1,2\}$,
$\{1,4\}$, and $\{2,4\}$, so $\{1,2,4\}$ is a third-order clique candidate.
It does not occur in $X_A$, whereas Section~\ref{sec:variability} already constructed the
system $X_B$ realizing this candidate.
Clique enumeration thus identifies the higher-order targets worth examining, while the
construction algorithm produces the complete records in which the target actually occurs.

\subsection{Integer realizability under finite resources}\label{subsec:boundary-ilp}

For given quotas $\boldsymbol c$, size distribution $(M_k)$, and forbidden set $\mathcal Z$,
define all legal complete block types
\begin{equation}
\mathcal T=\bigl\{B\subseteq V:
|B|\in\{k:M_k>0\},\ 
\binom B2\cap\mathcal Z=\varnothing\bigr\}.
\label{eq:allowed-types}
\end{equation}
Each type $B\in\mathcal T$ corresponds to a nonnegative integer variable $z_B$, the number of
occurrences of that type in the final system.

Under positive-association coverage requirements, the existence of a finite combinatorial
system can be written as the following integer model:
\begin{equation}
\begin{aligned}
\sum_{\substack{B\in\mathcal T\\i\in B}}z_B &=c_i,
&&i\in V,\\
\sum_{\substack{B\in\mathcal T\\|B|=k}}z_B &=M_k,
&&k=1,\ldots,n,\\
\sum_{\substack{B\in\mathcal T\\\{i,j\}\subseteq B}}z_B&\ge1,
&&\{i,j\}\in L_Q,\\
z_B&\in\mathbb Z_{\ge0},
&&B\in\mathcal T .
\end{aligned}
\label{eq:integer-feasibility}
\end{equation}
Every integer feasible solution corresponds to a multiset of blocks with the prescribed
quotas, size distribution, and positive-association coverage.
Conversely, every system satisfying these conditions yields an integer feasible solution.
If exact low-order conservation is required, the pairwise co-occurrence constraints and the
singleton counts are further fixed to $t_{ij}(X_0)$ and $s_i(X_0)$;
in tolerance mode, the constraint $E(z)\le\varepsilon$ computed from the $z_B$ is added to
the model.

Let $\mathcal Z_{\rm fin}$ denote the set of integer feasible solutions of
\eqref{eq:integer-feasibility}, and assume it is nonempty.
For $|S|\ge3$, define
\begin{equation}
U(S)=
\max_{z\in\mathcal Z_{\rm fin}}
\sum_{\substack{B\in\mathcal T\\S\subseteq B}}z_B .
\label{eq:max-higher-count}
\end{equation}

\begin{proposition}[Higher-order realizability under finite resources]
\label{prop:finite-realizability}
A prescribed higher-order set $S$ can occur in at least one combinatorial system satisfying
\eqref{eq:integer-feasibility} if and only if $U(S)\ge1$.
Moreover, $U(S)$ is its maximum realizable occurrence count under these constraints.
\end{proposition}
\begin{proof}
For any integer feasible solution $z$, the summand of \eqref{eq:max-higher-count} equals
exactly the number of complete records containing $S$ in the corresponding system.
Hence the existence of a system containing $S$ at least once is equivalent to the maximum of
this integer objective being at least one.
\end{proof}

Accordingly, define the sets of $h$-th order structures realizable under finite resources
\begin{equation}
\mathcal H_h^{\rm fin}
=\{S\in\binom Vh:U(S)\ge1\}.
\label{eq:finite-family}
\end{equation}
For a target $S\subseteq T$ we also have $U(T)\le U(S)$;
therefore, once the maximum occurrence count of a structure $S$ is zero, all its
higher-order extensions can be pruned directly during enumeration.

This model also connects with the residual subtraction of Section~\ref{sec:directed}.
If one complete record $B$ is fixed first, then the remaining element quotas become
$c_i-\mathbf1\{i\in B\}$ and the remaining size counts become
$M_k-\mathbf1\{|B|=k\}$.
The positive-association edges already covered by $B$ are also subtracted from the remaining
coverage task.
For a group of cumulatively fixed records $F$, these updates are written uniformly as
\begin{equation}
d_i=c_i-c_i(F),\qquad
M_k'=M_k-M_k(F),\qquad
\ell_{ij}=\max\{0,1-t_{ij}(F)\}.
\label{eq:feas-residual}
\end{equation}
Here $\ell_{ij}$ is the minimum co-occurrence count that the remaining system must realize.
A realizability decision can therefore fix candidate records one by one, update the resources
recursively, and solve the remaining integer feasibility problem at each layer.

\subsection{A resource-conflict example for a clique candidate}\label{subsec:boundary-counterexample}

Consider four elements $a,b,c,d$ and the combinatorial system
\begin{equation}
X_C=\bigl\{\{a,b,d\},\{a,c\},\{b,c\}\bigr\}.
\label{eq:resource-counterexample}
\end{equation}
Its size distribution is $M_3=1,M_2=2$, and its element quotas are
$(c_a,c_b,c_c,c_d)=(2,2,2,1)$.
The actual support graph contains $\{a,b\},\{a,c\},\{b,c\}$, so $\{a,b,c\}$ forms a triangle.

Now fix the quotas and the size distribution above and require every positive association of
the original support graph to occur at least once.
Element $d$ has only one occurrence, while the edges $\{a,d\}$ and $\{b,d\}$ both need to be
covered, so $a,b,d$ must jointly occupy the record containing $d$.
Since the system has only one $3$-element record, it must be $\{a,b,d\}$.
Hence
\begin{equation}
U(\{a,b,c\})=0.
\label{eq:resource-exclusion}
\end{equation}
This example shows that the association graph supplies combinatorial candidates, whereas the
finite record capacities, element quotas, and association coverage together determine whether
a candidate is actually realizable.

\subsection{From mathematical realizability to operational reachability}\label{subsec:boundary-reachability}

Fix an initial system $X_0$ and an operation set $\mathcal A$.
Take the complete systems satisfying the hard constraints as the vertices of a state graph,
and the operations of $\mathcal A$, applied one at a time, as the edges between vertices.
An operation may be a general legal $2$-switch, a zero-sum compensating trade as in
Section~\ref{sec:variability}, or may also include the freezing rules of
Section~\ref{sec:directed}.

Let $\mathcal R_{\mathcal A}^{\rm end}(X_0)$ denote the set of states reachable from $X_0$
by finitely many operation edges, with the hard constraints of $\Omega$ maintained along the
way; let $\mathcal R_{\mathcal A}^{\rm path}(X_0)$ denote the set of reachable states that
additionally require every intermediate state to satisfy a prescribed statistical tolerance,
association coverage, or freezing condition.
By definition,
\begin{equation}
\mathcal R_{\mathcal A}^{\rm path}(X_0)
\subseteq\mathcal R_{\mathcal A}^{\rm end}(X_0).
\label{eq:path-inclusion}
\end{equation}

This distinction can again be illustrated by $X_A$ and $X_B$.
In $X_A$ every nonzero element pair occurs exactly once.
Any nontrivial single $2$-switch changes the counts of some of these pairs, so it cannot serve
as an independent search step that preserves the complete second-order counts step by step.
After Section~\ref{sec:variability} combines two switches into one zero-sum operation, one
directly obtains an operation edge from $X_A$ to $X_B$ inside the equivalence class.
The composite operation thus enlarges the local transformations available under strict
low-order statistical constraints.

For a given operation rule, define the set of operationally reachable higher-order structures
\begin{equation}
\mathcal H_h^{\rm op}(X_0;\mathcal A)
=\bigcup_{\substack{X\in\mathcal R_{\mathcal A}^{\rm end}(X_0)\\
X\in\Omega^+(L_Q)}}
\{S\in\binom Vh:t_S(X)\ge1\}.
\label{eq:operation-family}
\end{equation}
Setting the operation endpoints under the same finite resources and association coverage as
\eqref{eq:integer-feasibility} gives
\begin{equation}
\boxed{\ 
\mathcal H_h^{\rm op}(X_0;\mathcal A)
\subseteq\mathcal H_h^{\rm fin}.
\ }
\label{eq:finite-operation-inclusion}
\end{equation}
This holds because every legal reachable endpoint is itself a feasible construction of the
integer model.
For problems of finite size one can search the state graph explicitly;
for larger problems, the compensating trades of Section~\ref{sec:variability}, the directed
construction of Section~\ref{sec:directed}, and integer feasibility checks together supply
structural witnesses.

This section thus completes the connection of the three layers:
clique enumeration gives graph-theoretic candidates for higher-order structures, the integer
model determines the actually realizable set under finite resources, and the state graph
characterizes the reachable part under the prescribed operation rule and initial conditions.
This boundary analysis provides a unified decision framework for the structural
transformations and directed constructions of the two preceding sections.

\end{document}
